\section{Numerical investigation}
\label{sec:numerical}




 In order to check numerically the  ideas discussed above we performed lattice
 QCD calculations in the quenched approximation at $\beta=6.0$ on $32^3\times 64 $
 lattices (lattice spacing $a=0.093$ fm). We used the non-perturbatively tuned clover fermion action with the
  clover coefficients computed by the Alpha  \mbox{collaboration \cite{Luscher:1996ug}.}

  We used a total of 500 configurations separated by 1000 updates each one consisting of four over-relaxation   and one
  heatbath sweeps.
 On each configuration we computed correlation functions from 6 randomly selected point sources.
  The pion  and nucleon masses in this setup were determined to be $601(1)$ MeV
  and $1411(4)$ MeV respectively. Conversion to physical energy units was
  performed used the Alpha collaboration scale setting
  for quenched QCD \cite{Necco:2001xg}.

  Our nucleon states were boosted up to a total momentum of $2.5\, $GeV (corresponding to the 6th lattice momentum).
  Inside  this momentum  range, the continuum dispersion relation for the nucleon was satisfied within the errors of the calculation,
   indicating small lattice artifacts of ${\cal O}(aP)$.
    In Fig. \ref{fig:disp} we plot the nucleon energy as a function of   momentum along with the continuum dispersion relation corresponding to our lattice nucleon zero momentum energy.

    \begin{figure}[t]
  \includegraphics[width=0.55\textwidth]{images/disperkey.pdf}
  \caption{Nucleon dispersion relation. Energies and momenta are in lattice units. The solid line is the continuum dispersion relation (not a fit) while the errorband is an indication of the statistical error of the lattice nucleon energies. }
  \label{fig:disp}
\end{figure}



 The computation of the matrix elements was performed using
  the methodology described in~\cite{Bouchard:2016heu} with an operator insertion given by Eq.~(\ref{Malpha}).
 Taking the time component of the current we can isolate  $\mathcal{M}_p(-z\cdot p, -z^2)$ which as discussed above is directly related to PDFs.

 Following~\cite{Bouchard:2016heu} we need to compute two types of correlation functions. The first is a regular nucleon two point function given by
%
 \begin{equation}
 C_P(t) = \langle \mathcal{N}_P(t)\overline{ \mathcal{N}}_P(0) \rangle\,  ,
 \end{equation}
%
where $ \mathcal{N}_P(t)$ is a helicity averaged, non-relativistic nucleon interpolating field with momentum $p$. The quark fields in $ \mathcal{N}_p(t)$ are smeared with a gauge  invariant Gaussian smearing. This choice of an interpolation field
is known to couple well to the nucleon ground state (see discussion in~\cite{Bouchard:2016heu}).
 The quark smearing width was optimized to give good  {overlap with the nucleon ground state within the}  range of momenta {in our calculation}.  The second correlator is given by
 \begin{equation}
 C^{\mathcal{O}^0(z)}_P(t) =\sum_\tau \langle \mathcal{N}_P(t) \mathcal{O}^0(z,\tau)\overline{ \mathcal{N}}_P(0) \rangle\, ,
 \end{equation}
 where
  \begin{equation}
 \mathcal{O}^0(z,t)= \overline\psi(0,t) \gamma^0 \tau_3 \hat{E}(0,z;A)\psi(z,t)\,,
 \end{equation}
 with $\tau_3$  being the flavor Pauli matrix.
 The proton momentum and the displacement of the quark fields were both taken along the $\hat z$ axis ($\vec z = z_3 \hat z$ and $\vec p = P \hat z$).
 We  define the effective matrix element as
 \begin{equation}
 \mathcal{M}_{\rm eff}(z_3 P, z_3^2;t) = \frac{C^{\mathcal{O}^0(z)}_P(t+1)}{C_P(t+1)} - \frac{C^{\mathcal{O}^0(z)}_P(t)}{C_P(t) }  \  .
 \end{equation}
 As it was shown in~\cite{Bouchard:2016heu}, our matrix element $\mathcal{J}$ can then be extracted at the large Euclidean time separation as
 \begin{equation}
\frac{\mathcal{J}(z_3P, z_3^2)}{2 E}= \lim_{t\rightarrow\infty}  \mathcal{M}_{\rm eff}(z_3 P, z_3^2;t) \  ,
 \end{equation}
 where $E$ is the energy of the nucleon.
      \begin{figure}[t]
  \includegraphics[width=0.48\textwidth]{images/plateau-p2z4.pdf} \hfill \includegraphics[width=0.48\textwidth]{images/plateau-p3z8.pdf}
  \caption{Typical fits used to extract the reduced matrix element. The upper panel corresponds to $p=2\pi/L \cdot 2$ and $z=4$ and the
  lower panel to  $p=2\pi/L \cdot 3$ and $z=8$, where momentum and position are in lattice units.
  }
  \label{fig:plateau}
\end{figure}
This method of extracting the matrix element, contrary to the traditional sequential source approach, allows for the computation of the matrix element using all source-sink separations for the nucleon creation and annihilation operators.

The  resulting effective matrix element has contaminations from excited states that scale as $e^{-t \Delta E }$,
where $t$ is the Euclidean time separation of the nucleon creation and annihilation operators,
and $\Delta E$ is the mass gap to the first excited state of the nucleon. Furthermore, it allows for the computation of  all nucleon matrix elements that correspond to different nucleon momentum spin polarization and flavor structure without additional computational cost.

As a  result,  the total computational cost of this approach is less than the equivalent cost of performing the calculations with the sequential source method, especially because in our approach we put emphasis on having as many nucleon momentum states as possible.  This approach has recently been successfully used for both single and multi-nucleon matrix element calculations~\cite{Berkowitz:2017gql,Tiburzi:2017iux,Shanahan:2017bgi}.

In order to normalize our lattice matrix elements we note that, for $z_3=0$, the matrix element
 $\mathcal{M}(z_3 P, z_3^2)$  corresponds to a  local vector (iso-vector) current, and therefore should  be equal to 1. However, on the lattice this is not the case due to lattice artifacts.
 Therefore we introduce a renormalization constant
 \begin{equation}
 Z_P = \frac{1}{ \left.\mathcal{J}(z_3 P, z_3^2)\right|_{z_3=0}}\, .
 \label{eq:renorm}
 \end{equation}
 The factor  $Z_P$ has to be independent from $P$.  However,  again due to lattice
  artifacts or potential fitting systematics,  this is not the case. For this reason,
   we renormalize  the matrix element for
   each momentum  with its own $Z_P$ factor taking this way advantage of maximal statistical correlations {to reduce statistical errors},  as well as the cancellation of lattice artifacts in the ratio.
  Therefore,  our  matrix element is extracted  using the ratio
  \begin{equation}
\mathcal{M}(z_3 P, z_3^2)= \lim_{t\rightarrow\infty}  \frac{\mathcal{M}_{\rm eff}(z_3 P, z_3^2;t)}{\left.\mathcal{M}_{\rm eff}(z_3 P, z_3^2;t)\right|_{z_3=0}}  \  .
 \end{equation}
 In order to determine the reduced matrix element $\mathfrak{M}(\nu,z_3^2)$ we introduce the double ratio
 \begin{align}
\mathfrak{M}(\nu, z_3^2)=&  \lim_{t\rightarrow\infty}
 \frac{\mathcal{M}_{\rm eff}(z_3 P, z_3^2;t)}{\left.\mathcal{M}_{\rm eff}(z_3 P, z_3^2;t)\right|_{z_3=0}} \nn
 &\times
\frac{\left.\mathcal{M}_{\rm eff}(z_3  P, z_3^2;t)\right|_{P=0,z_3=0}}
{\left.\mathcal{M}_{\rm eff}(z_3  P, z_3^2;t)\right|_{P=0}}  \ ,
\label{eq:redMatElem}
 \end{align}
 which takes care of the renormalization of the vector current according to Eq.~(\ref{eq:renorm}).
 In practice, the   infinite $t$ limit is obtained with a fit to a constant  for  a suitable choice of a fitting range. In all cases we studied,  the average $\chi^2$ per degree of freedom was ${\cal O}(1)$. Typical fits used to extract the reduced matrix element are presented in Fig.~\ref{fig:plateau}. All fits are performed with the full covariance matrix and the error bars are determined with  the jackknife method.

 We note here that that the reduced  matrix element defined in Eq.~(\ref{eq:redMatElem}) has a well defined continuum limit and no additional renormalization is required. This continuum limit is obtained at fixed $\nu$ and $z^2$ as well as at fixed quark mass.


In this calculation we used momenta up to $6\cdot 2\pi/L $ along the $z$-axis. This corresponds to a physical momentum of about 2.5\,GeV.
